\section{Supporting certificates and exact arithmetic checks}
\label{sec:support}

The analytic identities and counterexamples above do not depend on the
computations in this section. The numerical work serves two purposes:
to demonstrate that the complete theta source can satisfy the finite
boundary hypotheses, and to check implementations against exact models.
All reported numerical claims retain the scope and internal-review
status of their recorded calculations.

\subsection{A complete-theta application of the two-root theorem}

The interval calculation in \cite{note06,realcap} uses
\begin{equation}
 |\operatorname{Re}t|,|\operatorname{Im}t|\leq10^{-2},\qquad
 a=\frac{71}{2},\quad r=8,\quad R=10.
 \label{n:window}
\end{equation}
Both circular root counts are two. The enclosing data, rounded
outward for display, are
\begin{center}
\begin{tabular}{@{}lr@{}}
\toprule
Quantity & Certified bound\\
\midrule
Inner circular wall floor & $>8.36937\times10^{-7}$\\
Outer circular wall floor & $>7.57583\times10^{-7}$\\
Exterior coefficient modulus & $|A|<3.489078$\\
Phase product & $10^{-2}\sup|A|<0.034891$\\
Discriminant at each real time cap & $\Delta(\pm10^{-2})>144$\\
Complete source integrals & $4146$\\
\bottomrule
\end{tabular}
\end{center}
The two-root cap theorem therefore gives $Q=0$ on this window.
No RH assumption or tabulated zero location is an input.

The certificate covers entire arcs and the full complex-time rectangle,
not only sampled points. It uses 256 arcs on each circle, complete jets
of orders zero through fifteen at the time-zero anchors, and a
joint Taylor expansion of total degree seven. Complete moments pay the
order-eight remainder. All omitted source modes and exterior integration
tails are bounded. Zero-avoiding convex disks around the boundary image
arcs justify the polygonal winding counts. The coefficient estimate
retains the coupled numerator $H_z+(3/10)H$ before taking absolute
values. If $d_k$ and $n_k$ are the full-arc lower and upper bounds for
the denominator and this numerator, respectively, the calculation uses
\[
 |A|\leq\frac{4R}{(R-r)^2}\frac1{256}\sum_{k=0}^{255}\frac{n_k}{d_k}.
\]
At each real time cap, the certified signs at spatial points
$28,29,41,43$ are $+,-,-,+$. The intermediate value theorem and the
total count of two put one simple root in each of $(28,29)$ and
$(41,43)$. Their gap exceeds twelve, proving the displayed
discriminant bound. Thus the cap signs have an independent real-source
payment.

This is a separated-pair example. Its large cap gap does not address an
unknown shrinking gap near a hypothetical endpoint. The same boundary
and phase procedure also admits the exact quadratic collision
$H=(z-71/2)^2-2t$, returning charge one from its opposite cap signs.
It therefore does not silently assume that events are absent.

\subsection{Earlier direct comparisons}

The earlier finite experiments \cite{note01,note03} compared specified
boundary and direct joint-exclusion procedures using complete source
data. Their certified margins and acquisition counts were
\begin{center}
\begin{tabular}{@{}lccc@{}}
\toprule
Experiment & Boundary margin & Direct margin & Source integrals\\
\midrule
One-zero ellipsoid & $>7.495\times10^{-5}$ & $>7.993\times10^{-5}$ & $8$ versus $6$\\
Two-zero thin box & $>3.5683\times10^{-7}$ & $>3.9348\times10^{-7}$ & $1032$ each\\
\bottomrule
\end{tabular}
\end{center}
The thin box has $|\operatorname{Re}t|,|\operatorname{Im}t|\leq0.01$,
$28\leq\operatorname{Re}z\leq43$, and
$|\operatorname{Im}z|\leq0.01$. Both methods used the same source
acquisition and 2056 boundary vertices. The direct procedure had the
stronger certified margin in these declared implementations. This is
neither a universal complexity ordering nor evidence against other
boundary geometries. The analytic reference-transfer result in
Appendix~\ref{app:reference} explains why an already certified
holomorphic reference does not by itself confer an improved uniform
margin on the boundary method.

\subsection{Exact controls and reproducibility}

The standard-library rational calculation \cite{controls} checks the
Hermite discriminants for $2\leq m\leq8$, five finite clusters with
nonzero exterior forcing, the cubic expression for $T$, and quadratic
source response. It also verifies the explicit boundary-paid affine
homotopy example from the signed-response section. These checks agree
with the proofs, including the negative-time signs of the many-root
forcing. They are finite algebra checks, not a proof for all $m$.

The earlier piecewise affine boundary-degree controls \cite{note02}
return degrees $+1,-1,+1,0$ for, respectively, a complex-oriented
linear map, an orientation-reversing real map, an enclosed quadratic
heat event, and a quadratic heat field with its event outside the box.
The negative-degree example deliberately fails holomorphicity. It
checks the orientation convention rather than holomorphic positivity.

The repository records supply scripts, arithmetic versions, parameter
choices, source-tail and interpolation budgets, and small execution
receipts. The theta runs use \texttt{python-flint 0.9.0}; the exact
algebra checks require only Python's standard library. Source-array
hashes bind the recorded computations but do not replace their
analytic error estimates. No additional theta numerical region was
acquired for this manuscript.

\section{What remains to obtain endpoint exclusion}
\label{sec:open}

The geometric difficulty at an isolated hypothetical double has been
resolved in a pointwise quantitative sense: a positive radius, collar,
time width, pair of root counts, and phase margin follow from the
nonzero second jet and complete moments. The resulting window correctly
has charge one when a double is assumed. A new exclusion must therefore
use information beyond those universal local properties.

\paragraph{A sign for the unchanged theta source.}
The signed response and the canonical-tail mixture retain correlations
that absolute moment bounds discard. Their remaining task is to prove,
with a declared reference and the entire residual included, either a
strictly positive incoming discriminant or a strictly negative event
location. For a proper two-root window with a positive outgoing cap,
the first gives charge zero. Alternatively, an independently located
negative-time event uses the unique event allowed by the phase theorem
and excludes another event at zero. Neither required theta-specific
inequality has been established here. A strict improvement over another
enclosure can still leave zero inside the resulting interval.

\paragraph{Height coverage.}
The windows exist separately at every finite hypothetical double
height. Their sizes depend on the local curvature; the displayed
construction can shrink with height. The source sign, error terms,
tail cutoff and event-location margin must therefore be controlled with
the same dependence. A calculation on one fixed interval, or an
unquantified limiting statement about the head cutoff, cannot discharge
the universal quantifier over endpoint heights. No common negative
all-real time slice is being proposed.

\paragraph{Every multiplicity.}
The all-root identity $D'=8T+BD$ retains the whole cluster and its
exterior contribution. For $m>2$, the forcing $T$ is no longer constant.
The cubic control shows that nonreal time events survive even when
$B=0$; the quartic control has charge six despite equal positive real
cap signs. Applying the two-root theorem to an arbitrarily selected
pair at a higher birth is invalid, because the other birth roots
destroy the required isolating collar. The cubic real-event monotonicity
provides a possible conditional event-location argument, but no
complete-source location estimate is supplied. Higher multiplicities
still require an exclusion compatible with their full discriminant.

The finite-moment control constructed above has exactly its stated
preservation properties; it does not realize the actual theta moment
vector. Broader collision controls in the project have different
failures of the theta premises: Euler-factor controls change arithmetic
data; finite square-source and finite-prime controls lack the complete
normalized coefficient/reflection conjunction; reflected or smooth
positive-source controls lack the canonical square-source expansion;
the close positive endpoint controls in \cite{bw10control} fail ordinary
Dirichlet-series arithmetic. Their preservation lists must not be
combined. They rule out arguments using only their shared premises,
not estimates specific to the unchanged complete theta source.

In the task register, the hypothetical-double window is complete in
its stated parameterized scope (ECE-01). The distinguishing sign
(ECE-02), its all-height use (the signed part of ECE-04), and
higher-multiplicity exclusion (the remaining part of ECE-05) are open.
Endpoint exclusion and RH-conditional simplicity (ECE-03 and ECE-06)
remain contingent on those results. The manuscript is an analytic
consolidation, not an upgrade of those statuses.

\section*{Acknowledgements and research provenance}
\addcontentsline{toc}{section}{Acknowledgements and research provenance}
This manuscript was drafted for Edward Baker with substantial assistance
from OpenAI's Codex language-model system. The manuscript drafting and
internal review session recorded model \texttt{gpt-6-astra}, reasoning
effort \texttt{ultra}. The underlying 21 October 2026 notes record their
own configurations, including earlier \texttt{max} sessions; those
provenance distinctions are retained in the source records. Separate
agents of the same model assisted in drafting and checking sections.
These checks do not constitute independent specialist review or formal
verification. The author designation above identifies the intended
recipient of the draft and does not imply that he has approved every
claim. The investigation's notes, reviews, progress register and chat
archive preserve the detailed research history; no manuscript snapshot
series is created.
